\Article{34}
Es würde von unserem gegenwärtigen Vorhaben wegführen, von den hier aufgestellten allgemeinen Grundsätzen zu einzelnen Erscheinungen überzugehen. Das gilt besonders, weil diese Grundsätze ihrem Wesen nach mit der Theorie übereinstimmen, mit der der berühmte \textsc{Laplace} bereits sehr viele Erscheinungen des Flüssigkeitsgleichgewichts mit gleichem Geschick und Erfolg erklärt hat. Es bleibt gewiss ein weites Feld, das eine reiche neue Ernte verspricht. Sie sei jedoch künftigen Arbeiten vorbehalten. Angebracht ist es dagegen, einige Bemerkungen hinzuzufügen, die entweder neues Licht auf den Gegenstand werfen oder eine falsche Deutung verhindern können.

I. Unsere Theorie erhebt keinen Anspruch darauf, die Gleichgewichtsform mathematisch genau zu bestimmen. Sie begnügt sich damit, eine Gestalt zu bestimmen, von der die wahre Gleichgewichtsform nicht um einen merklichen Betrag abweichen kann. Es wäre ein Irrtum, dies einer Unvollkommenheit der Theorie zuzuschreiben. Sie hat vollständig geleistet, was sich leisten lässt, solange das Gesetz der molekularen Anziehung unbekannt ist. Im Gleichgewichtszustand muss die Funktion $\Omega$ genau ein Maximum sein und folglich die Funktion
\[
\frac{2\pi cs\psi0}g-\frac\Omega{gc}
\]
ein Minimum. Wegen der Beschaffenheit der molekularen Anziehung ist diese Funktion zwar nicht genau gleich der Funktion $W$, weicht aber nur unmerklich von ihr ab. Die Gestalt, für die $W$ ein Minimum wird, ist also nicht die genaue Gleichgewichtsform. Der Unterschied muss jedoch unmerklich sein, insofern jede merkliche Änderung dieser Gestalt einen merklich größeren Wert der Funktion $W$ ergeben würde. Damit ist offenbar ein merklicher Unterschied in der Krümmung der Oberfläche möglich, sofern er sich auf einen unmerklichen Teil der Oberfläche beschränkt. Bei der genauen Gleichgewichtsform darf man daher den oben mit $A$ bezeichneten konstanten Winkel nur in einem nicht messbaren Abstand vom Gefäß als Neigung der Flüssigkeitsoberfläche gegen die Gefäßwand auffassen, und nicht unmittelbar an der Berührungsstelle. Anders gesagt, an der Grenze des Bereichs, in dem die Anziehung des Gefäßes merklich ist, stimmt die Neigung bis auf einen unmerklichen Unterschied mit dem Wert von $A$ überein, wie der berühmte \textsc{Laplace} bereits richtig bemerkt hat.

\SourceLeaf{79}II. Man muss die Gleichgewichtsform sorgfältig von der Gestalt im Ruhezustand unterscheiden. Befindet sich eine Flüssigkeit im Gleichgewicht, so wird sie gewiss darin verharren. Weicht ihre Gestalt aber etwas von der Gleichgewichtsform ab, so kann die Flüssigkeit dennoch in Ruhe bleiben oder, wenn sie sich bewegt, schon vor Erreichen des Gleichgewichts zur Ruhe kommen. Ebenso ist beispielsweise ein Würfel, den man einfach auf eine horizontale Ebene setzt, im Gleichgewicht. Er kann aber auch auf einer geneigten Ebene ruhen, wenn Reibung die Bewegung verhindert. Eine Flüssigkeit in einem Zustand, für den $W$ einen kleinsten Wert hat, wird also gewiss ruhen. Befindet sie sich dagegen in einem anderen Zustand, etwa wenn $W$ abnehmen kann, so geht sie aus diesem Zustand nur insoweit in den Gleichgewichtszustand über, als Reibung sie daran nicht hindert. In dieser Hinsicht unterscheiden sich die beiden Gleichgewichtsbedingungen wesentlich. Die erste Grundgleichung (Artikel 28) hängt nicht davon ab, ob sich die Grenze $P$ verschieben kann. Sie ist also auch dann für die Bedingung des Minimums notwendig, wenn man die Grenze als fest voraussetzt. Soweit die Flüssigkeit vollkommen frei fließen kann, sodass ein Teil frei über einen anderen gleiten kann, sobald auch nur die kleinste Kraft eine Bewegung verlangt, wird sie sich dieser Bedingung notwendig anpassen. Beim zweiten Grundsatz (Artikel 29) liegen die Dinge ganz anders. Er hängt wesentlich von der vollkommenen Beweglichkeit der Grenze $P$ auf der Gefäßoberfläche ab. Die Bedingung, dass der Wert von $W$ ein Minimum ist, verlangt durchaus die Gleichung $i=A$. Hat sich aber die Flüssigkeitsoberfläche dem ersten Grundsatz angepasst, ohne dass der Winkel $i$ schon seinen Normalwert und damit $W$ seinen absoluten Kleinstwert erreicht hat, so ist der Übergang zum vollkommenen Gleichgewicht nur durch eine Verschiebung der Grenze $P$ möglich, also durch eine Bewegung der Flüssigkeit an der Berührungsstelle mit dem Gefäß. Einer solchen Bewegung kann die Reibung entgegenstehen. Daraus wird verständlich, warum wir bei Versuchen mit denselben Körpern so große Unterschiede im Wert des Winkels $i$ antreffen. Ebenso wird sich im Fall $\beta\beta>\alpha\alpha$ eine Flüssigkeit in einem Gefäß, dessen Wände bereits benetzt sind, nach dem Gleichgewichtsgesetz einstellen, dem zufolge für den Hauptteil der Flüssigkeit $i=180^0$ gelten muss. Sind dagegen die Gefäßwände außerhalb der Flüssigkeit noch trocken, so kann eine Flüssigkeit, die von einem Zustand außerhalb des Gleichgewichts ausgeht und sich auf den trockenen Gefäßwänden ausbreitet, bereits zur Ruhe kommen, bevor der Winkel $i$ den Wert $180^0$ erreicht hat. Das erklärt zugleich, warum die Kapillarerscheinungen von Flüssigkeiten, die einer Benetzung nicht entgegenstehen, in trockenen Röhren so große Unregelmäßigkeiten zeigen und der Anstieg sehr oft viel geringer ist als in schon benetzten Röhren, in denen wir stets die schönste Übereinstimmung mit der Theorie sehen.

\expandafter\def\csname HistoricalNoteText7\endcsname{Es ist zu beachten, dass die Größe, die der berühmte \textsc{Laplace} mit $H$ bezeichnet, mit unserem $\pi c\theta0$ übereinstimmt. Sein $a$ bezeichnet daher dasselbe, was in unserer Schreibweise $\frac{g}{\pi c\theta0}$ oder $\frac1{2\alpha\alpha}$ ist.}%
\SourceLeaf{80}III. Das Verhältnis der Konstanten $\alpha,\beta$ lässt sich aus den Erscheinungen nicht ableiten, wenn $\beta$ größer als $\alpha$ ist. Denn in diesem Fall unterscheidet sich die Gestalt derselben Flüssigkeit in Gefäßen gleicher Form aus verschiedenen Materialien nur hinsichtlich des nicht messbaren Films, der das Gefäß benetzt. Ist dagegen $\beta$ kleiner als $\alpha$, so lässt sich das Verhältnis dieser Konstanten mithilfe des Winkels $i$ zwar bestimmen, doch ist aus den eben genannten Gründen kaum große Genauigkeit zu erwarten. Für Quecksilber in Glasgefäßen setzte der berühmte \textsc{Laplace} den Winkel $i=43^0\,12'$ an.

Die Konstante $\alpha$ lässt sich viel genauer bestimmen, besonders wenn wir Gefäße verwenden können, die sich benetzen lassen. Für Wasser bei einer Temperatur von $8{,}5$ Grad auf dem hundertteiligen Thermometer ist nach den vom berühmten \textsc{Laplace} angeführten Versuchen anzusetzen\HistoricalNote{7}
\[
\begin{aligned}
\alpha\alpha&=7{,}5675\ \text{Quadratmillimeter, oder}\\
\alpha&=2{,}7509\ \text{Millimeter}
\end{aligned}
\]
Für Weingeist mit dem spezifischen Gewicht $=0{,}81961$ bei derselben Temperatur gilt
\[
\begin{aligned}
\alpha\alpha&=3{,}0441\ \text{Quadratmillimeter, oder}\\
\alpha&=1{,}7447\ \text{Millimeter}
\end{aligned}
\]
Für Terpentinöl bei einer Temperatur von $8$ Grad gilt
\[
\begin{aligned}
\alpha\alpha&=3{,}305\ \text{Quadratmillimeter, oder}\\
\alpha&=1{,}818\ \text{Millimeter}
\end{aligned}
\]
Für Quecksilber bei einer Temperatur von $10$ Grad dürfen wir, bis neue Versuche größere Genauigkeit liefern, ansetzen
\[
\begin{aligned}
\alpha\alpha&=3{,}25\ \text{Quadratmillimeter, oder}\\
\alpha&=1{,}803\ \text{Millimeter}
\end{aligned}
\]
Im Übrigen ist es wahrscheinlich, dass die Temperatur den Wert der Konstanten $\alpha\alpha$ nur insoweit beeinflusst, als die Dichte von ihr abhängt. Nach dieser Annahme ist der Wert von $\alpha\alpha$ also proportional zur Dichte.

\SourceLeaf{81}Diese Werte wurden aus dem Anstieg oder der Absenkung von Flüssigkeiten in Kapillarröhren erschlossen. Es ist jedoch sehr schwierig, die Durchmesser dieser Röhren genau zu messen, und noch schwieriger, sich zu vergewissern, dass ihre Querschnitte kreisförmig sind. Viel größere Genauigkeit versprechen Versuche über Durchmesser und Volumina großer Quecksilbertropfen, die auf einer horizontalen Platte oder auf einer Platte von sehr geringer bekannter Krümmung ruhen. Solche Versuche haben die Physiker \textsc{Segner} und \textsc{Gay-Lussac} bereits angestellt. Dasselbe gilt bei Flüssigkeiten, die Glasgefäße benetzen, für Versuche über die Abmessungen großer Luftblasen in Gefäßen, die oben mit einem benetzten Deckel verschlossen sind, der eben und horizontal oder nach einem bekannten Radius leicht gekrümmt ist. Wir laden die Physiker ein, diese Versuche anzustellen.

IV. Um den dieser Abhandlung gesetzten Rahmen nicht zu überschreiten, mussten wir die Anwendung unserer allgemeinen Grundsätze hier auf den einfachsten Fall beschränken, in dem eine einzige Flüssigkeit in einem starren Gefäß betrachtet wird. Es steht jedoch nichts im Wege, der Theorie ihre volle Allgemeinheit zu geben, sodass sie auch das Problem mehrerer Flüssigkeiten im selben Gefäß und den Fall umfasst, dass außerdem ganz oder teilweise freie starre Körper in die Flüssigkeit eintauchen. Eine ausführlichere Darstellung dieser Fragen müssen wir uns aber für eine andere Gelegenheit vorbehalten.
